\section*{Geometry, constraints, and strong fields}
Gauge covariance identifies compatible transport. Gravity asks for more: the geometry, constraints, and evolution must be specified together. A propagation rule, a canonical constraint algebra, and an evolved matter system are different mathematical objects, with different hypotheses.

This account proceeds from weighted transport geometry and reduced Gowdy evolution to finite linearized ADM constraints. The nonlinear obstruction explains why a nonzero derivative cannot satisfy the ordinary Leibniz rule on a complete finite pointwise algebra. Matrix moment maps, homological transfer, and spectral representations provide alternatives in their stated settings.

A real Lorentzian action supplies a classical gauge-theoretic parent. Boundary constructions and local second-class reduction retain their regularity and nondegeneracy conditions. The spherical sector makes mass, flux, and horizon diagnostics explicit, allowing exact identities to be compared with bounded numerical experiments.

The results do not amount to a complete continuum or quantum theory of gravity. A manufactured electrovac trajectory is not an existence theorem for nonlinear collapse. A scalar trapped-sphere diagnostic is not a continuum formation theorem without stability, convergence, and event control. These distinctions are part of the mathematics rather than qualifications added after it.

The shared resolution language continues in the arithmetic and polyphase constructions that follow. It makes the state spaces, comparison maps, and residual obstructions explicit without equating their physical interpretations.

\clearpage
\section*{Reading routes and conventions}
\addcontentsline{toc}{section}{Reading routes and conventions}
For the geometric route, begin with transport and local frames, then follow the Gowdy model and the finite linearized ADM system. The nonlinear obstruction explains why a product rule that is harmless in a continuum calculation becomes a substantive issue on a finite pointwise algebra.

For the constrained-system route, begin with the canonical bracket and the exact finite moment map. Continue through homological transfer, resolved spectral evaluation, the classical BV parent, and the boundary and local reduction conditions. A field-dependent change of ghost variables and a canonical transformation of the extended phase space must be distinguished. The latter carries the appropriate momentum terms with it.

For the strong field route, begin with the spherical mass identity. The charged horizon calculations and the scalar mass-flux compatibility give concrete statements that can be checked independently. The numerical records then test selected trajectories and initial data. A manufactured gauge trajectory of a fixed spacetime and an evolved scalar-collapse model answer different questions.

\begin{longtable}{@{}L{.25\linewidth}L{.66\linewidth}@{}}
\toprule
Convention & Meaning\\\midrule\endhead
Canonical brackets & Each sector declares its variables and sign. In standard coordinates, $\{F,G\}=\sum_a(F_{q_a}G_{p_a}-F_{p_a}G_{q_a})$, so $\dot F=\{F,H\}$.\\
Geometry and signature & Lorentzian signature, curvature orientation, density weights, lapse, shift, and tetrad conventions are stated where they enter. Local transport alone does not select an Einstein action.\\
Gauge and reduction & Gauge choice, constraint enforcement, and removal of second-class variables are separate operations. An invertible gauge-fixing or Dirac matrix defines the regular chart used in the corresponding argument.\\
Spectral resolution & $K$ denotes retained degree and $N$ the sampling resolution in the resolved constructions. Conditions such as $16K<N$ are sufficient degree budgets in a declared representation. They are not invariant under every nonlinear change of variables.\\
Locality & A finite-range stencil, a full spectral evaluation, a matrix commutator, and a radial prefix constraint solve have different dependence on the data. The implementation is named before its locality is assessed.\\
Time stepping & The conservation proof for a projected discrete-gradient step concerns an autonomous Hamiltonian and an exactly solved implicit step. Approximate solves retain a residual that must be measured.\\
Spherical units & Geometrized units are used in the spherical examples. The mass, charge, radial coordinate, and slicing branch are defined before horizon or flux formulas are used.\\
Evidence & Ordinary proofs, inherited executable records, newly run bounded computations, and archived formal declarations retain their own scope. Numerical plots do not substitute for an existence or convergence proof.\\
\bottomrule
\end{longtable}

\clearpage
\begingroup
\setlength{\parskip}{0pt}
\setlength{\cftbeforesecskip}{4pt}
\setlength{\cftbeforepartskip}{10pt}

\endgroup
\clearpage

